\chapter{Low-Latency Inference}\label{ch:ml:low-latency-inference}

A gradient-boosted model of 300 trees answers in 265 microseconds through its library's Python interface and in 2.6
microseconds as generated C++, measured on the same laptop; the model did not change, and the two answers are equal to the
last bit. Most of the first figure is not the model: it is the interpreter, the wrapper, the conversion of one row into an
array and back. This chapter is about taking models out of that environment: compiling trees into code, shrinking small
networks to eight-bit integers, deciding when to batch, and doing it in C++20 and Rust with the exact results of the Python
reference. It rests on Book~13's method for measuring latency, and its timings are one laptop's, labelled as such; what the
tests check is parity, and what the chapter concludes is about orders of magnitude.

\section{The inference budget}

\begin{definition}[Inference latency]\label{def:ml:low-latency-inference:latency}
\emph{Inference latency}\index{inference latency} is the time from the moment a model's inputs are available to the moment
its output is, measured per call and reported as a distribution (median and tail), within the strategy's latency budget
(Book~13, chapter~1).
\end{definition}

A market-making or execution model decides inside a budget of a few microseconds to a few hundred, shared with feature
computation (chapter~24), risk checks and the network. Inference must fit its share at the tail, not the median: the decision
that arrives late is the one taken when the market moved. The chapter's two models are those of \texttt{firm.mlinfer}'s
fixture, fitted to a synthetic problem with 16 features: a LightGBM forest (chapter~5) of 300 trees with 15 leaves each, 4\,200
splits in all, and a \omterm{def:ml:neural-networks-for-noisy-tabular-data:mlp}{multilayer perceptron} (chapter~7) of $16\to32\to16\to1$ units. On 20\,000 new observations the forest's
rank information coefficient is 0.224 and the network's 0.182.

\section{Compiling trees}

\begin{definition}[Model compilation]\label{def:ml:low-latency-inference:compile}
\emph{Model compilation}\index{model compilation} translates a trained model into code or data structures specialised for
inference on a target (generated source, a flat array layout, a hardware description), removing the training framework from
the prediction path (Asadi, Lin and de Vries, 2014; Lucchese and co-authors, 2015, for ranking forests).
\end{definition}

The forest is exported two ways. As flat arrays of feature indices, thresholds and child indices, one entry per node, which a
loop walks (\cref{lst:ml:inf:loop}): the same loop runs in Python, C++20 and Rust. And as generated C++ with one nested
\texttt{if} per split and the thresholds as literals (\cref{lst:ml:inf:gen}), 17\,107 lines that the compiler turns into
straight branches. Both reproduce LightGBM's predictions exactly on the 200 shared test vectors, because they visit the same
leaves and add the leaf values in the same order in double precision; the thresholds survive the trip through text because
they are written with the shortest representation that reads back to the same double.

\Cref{fig:ml:inf:latency} shows the measured latencies. LightGBM's scikit-learn interface takes 265~µs per single-row call at
the median and 544~µs at the 99th percentile; the same flat loop in Python is no faster (300~µs), because the interpreter's
cost per node dominates. In C++ the loop takes 6.2~µs and the generated branches 2.6~µs (5.7~µs at the 99th percentile), a
hundred times less than the library call at both the median and the tail; the Rust loop takes 7.8~µs. The branches win over
the loop because the thresholds and child addresses are in the instruction stream rather than in memory, and the compiler
can lay out the likely path; with 300 trees and 4\,200 splits the whole forest still fits in the caches (Book~13, chapter~3).

\begin{omfigure}
\begin{tikzpicture}
\begin{semilogxaxis}[width=11cm,height=7.5cm,xbar,bar width=5pt,bar shift=0pt,xmin=10,xmax=2e6,
  xlabel={latency per call (nanoseconds)},ytick={0,1,2,3,4,5,6,7,8,9},
  yticklabels from table={figdata/ml/26-low-latency-inference/measured_latency.csv}{path},y dir=reverse,
  tick label style={font=\small},label style={font=\small},grid=major,grid style={gray!25},area legend,
  legend style={font=\footnotesize,at={(0.5,-0.2)},anchor=north,draw=none,fill=none},legend columns=2,
  table/col sep=comma,enlarge y limits=0.06,log origin=infty]
\addplot[fill=omThm,draw=none,bar shift=2.6pt] table[x=p99_ns,y=i]{figdata/ml/26-low-latency-inference/measured_latency.csv};
\addplot[fill=omDef,draw=none,bar shift=-2.6pt] table[x=median_ns,y=i]{figdata/ml/26-low-latency-inference/measured_latency.csv};
\legend{99th percentile,median}
\end{semilogxaxis}
\end{tikzpicture}

\omcaption{Single-call latency of the same two models by execution path; 200\,000 calls in C++ and Rust, 5\,000 in Python (300
for the pure-Python forest), timed one by one. Measured on a laptop (Intel Core Ultra 7 155H) under WSL2 at low priority, no
isolated cores; other jobs may have run. Data: \texttt{bench\_infer.py}, \texttt{measured\_latency.csv}.}\label{fig:ml:inf:latency}
\end{omfigure}

\section{Small networks: layout, fusion and quantisation}

\begin{definition}[Operator fusion, quantisation, post-training quantisation, quantisation-aware training]\label{def:ml:low-latency-inference:quant}
\emph{Operator fusion}\index{operator fusion} computes several consecutive operations (a matrix product, a bias, an
activation, a rescaling) in one pass over the data. \emph{Quantisation}\index{quantisation} represents weights and
activations as small integers with scale factors, so that inference runs in integer arithmetic (Jacob and co-authors, 2018).
\emph{Post-training quantisation}\index{post-training quantisation} derives the integers and scales from a trained float model
and a calibration sample; \emph{quantisation-aware training}\index{quantisation-aware training} fine-tunes the model with the
rounding simulated in the forward pass, so that it learns weights that survive it (Nagel and co-authors, 2021).
\end{definition}

The chapter's network is quantised symmetrically, one scale per tensor: weights and activations to integers in $[-127, 127]$,
biases to 32-bit integers at the product of the input and weight scales, and between layers a fixed-point multiplier and shift
that bring the 32-bit accumulator back to eight bits, rounding half away from zero (\cref{lst:ml:inf:mlp}). Each layer is one
fused pass: products, bias, rescaling and ReLU. The integer arithmetic makes parity exact by construction: the NumPy
reference, the C++ kernel and the Rust kernel agree on every one of the 200 test vectors to the bit. The eight-bit weights
take 1\,236 bytes against 4\,356 in float32, and one call takes 0.47~µs in C++ and 0.58~µs in Rust.

\begin{table}[htbp]
\centering\small
\begin{tabular}{@{}lccc@{}}
\toprule
& \multicolumn{3}{c}{bits per weight and activation}\\
\cmidrule(l){2-4}
rank IC on 20\,000 new observations & 8 & 6 & 4\\
\midrule
quantised after training & 0.182 & 0.181 & 0.143\\
quantisation-aware \omterm{def:ml:representation-learning:ssl}{fine-tuning}, then quantised & 0.179 & 0.176 & 0.164\\
\bottomrule
\end{tabular}
\caption{The network's information coefficient after \omterm{def:ml:low-latency-inference:quant}{quantisation} (float: 0.182; the forest: 0.224). Data:
\texttt{ml\_infer.accuracy}.}\label{tab:ml:inf:quant}
\end{table}

\Cref{tab:ml:inf:quant} holds the second half of the chapter's named result. At eight bits, \omterm{def:ml:low-latency-inference:quant}{quantisation} after training costs
nothing measurable (0.1817 against 0.1822; the integer model's outputs differ from the float model's by 4.5\% of their standard
deviation, and the ranking barely moves); quantisation-aware \omterm{def:ml:representation-learning:ssl}{fine-tuning}, which spends ten \omterm{def:ml:neural-networks-for-noisy-tabular-data:backprop}{epochs} adapting to the rounding,
loses a little (0.179), because the \omterm{def:ml:representation-learning:ssl}{fine-tuning} itself moves the model. At six bits the story is the same. At four, \omterm{def:ml:low-latency-inference:quant}{post-training
quantisation} loses a fifth of the signal (0.143) and \omterm{def:ml:low-latency-inference:quant}{quantisation-aware training} recovers more than half of the loss (0.164).
Eight bits is free for a network this small; fewer bits are where training for the rounding pays.

\section{Batching against latency}

\begin{definition}[Dynamic batching]\label{def:ml:low-latency-inference:batch}
\emph{Dynamic batching}\index{dynamic batching} groups requests that arrive close together into one call of the model, trading
the waiting time of the first request for a lower cost per request.
\end{definition}

For a library whose per-call overhead dwarfs its per-row cost, batching is a large lever (\cref{fig:ml:inf:batch}): LightGBM's
call takes 273~µs for one row and 3.1~ms for 1\,024, so the cost per row falls from 273~µs to 3.0~µs. A research job scoring a
universe each second should batch. A trading decision on one instrument cannot wait for others' requests to arrive, and a
compiled model already costs a few microseconds per row: there, batching buys throughput the strategy does not need at the
price of latency it cannot afford. When batching is used, latency is measured from each request's arrival, not from the
batch's start, or the waiting time vanishes from the statistics (coordinated omission, Book~13, chapter~5).

\begin{omfigure}
\begin{tikzpicture}
\begin{loglogaxis}[width=10.5cm,height=6cm,xlabel={rows per call},ylabel={microseconds},tick label style={font=\small},
  label style={font=\small},xmin=0.8,xmax=1300,grid=major,grid style={gray!25},legend cell align=left,
  legend style={font=\footnotesize,at={(0.02,0.02)},anchor=south west,fill=white,draw=none},table/col sep=comma]
\addplot[omThm,very thick,mark=*] table[x=batch,y=call_us]{figdata/ml/26-low-latency-inference/measured_batch.csv};
\addplot[omDef,very thick,mark=square*] table[x=batch,y=per_row_us]{figdata/ml/26-low-latency-inference/measured_batch.csv};
\legend{per call,per row}
\end{loglogaxis}
\end{tikzpicture}

\omcaption{LightGBM's prediction latency (median) against batch size through its Python interface, per call and per row.
Same machine and conditions as \cref{fig:ml:inf:latency}. Data: \texttt{bench\_infer.py},
\texttt{measured\_batch.csv}.}\label{fig:ml:inf:batch}
\end{omfigure}

\section{Inference in C++, in Rust and on programmable hardware}

The kernels load their model once, allocating, and then predict without touching the heap: fixed-size buffers on the stack, no
exceptions, no virtual calls, the model's arrays read-only and shared. The Rust version (\cref{lst:ml:inf:rust}) is the same
algorithm with bounds checks, which cost it about a quarter against C++ here; both could remove them with unchecked indexing
once the model's arrays are validated at load time. Beyond the CPU, programmable hardware runs trees and small networks in a
fixed number of clock cycles, laid out as a circuit: Duarte and co-authors (2018) compiled small networks for FPGAs in the
trigger systems of particle physics, where decisions have fixed latency budgets of the same kind. The same parity discipline
applies: a hardware model is validated against the Python reference on shared vectors, bit for bit where the arithmetic is
integer.

\begin{method}[Putting a model in the fast path]\label{met:ml:inf:method}
\begin{enumerate}
\item Measure the library call's median and tail on the target machine, one request at a time, before changing anything.
\item Export the model to a portable form and compile it (trees to code or flat arrays; networks to fused integer kernels),
with no allocation on the hot path.
\item Test parity on shared vectors: exact for trees and integer networks, within stated tolerances for floating-point
networks.
\item Quantise after training and measure the loss of signal; use \omterm{def:ml:low-latency-inference:quant}{quantisation-aware training} only when the loss matters.
\item Register the compiled artefact with its parity report next to the model it came from (chapter~25).
\end{enumerate}
\end{method}

\section{Tutorial: two microseconds for a forest}\label{tut:ml:low-latency-inference:tut}

\begin{tutorial}
\textbf{Goal.} Export a forest and a network, compile them to C++20 and Rust, check parity, quantise, and measure.
\textbf{End state:} \cref{fig:ml:inf:latency,fig:ml:inf:batch}, \cref{tab:ml:inf:quant}.
\begin{tutsteps}
\item \textbf{Trees as generated code.}
\omcode{code/firm/mlinfer/firm_mlinfer.py}{84}{103}{Generating nested C++ branches from the exported forest.}\label{lst:ml:inf:gen}
\item \textbf{The flat forest in C++20.}
\omcode{code/firm/mlinfer/cpp/mlinfer.hpp}{20}{27}{Walking the flat node arrays.}\label{lst:ml:inf:loop}
\item \textbf{The int8 network in C++20.}
\omcode{code/firm/mlinfer/cpp/mlinfer.hpp}{66}{94}{Integer inference with fixed-point requantisation.}\label{lst:ml:inf:mlp}
\item \textbf{The flat forest in Rust.}
\omcode{code/firm/mlinfer/rust/src/lib.rs}{15}{26}{The same loop in Rust.}\label{lst:ml:inf:rust}
\item \textbf{Run} \texttt{make\_mlinfer\_fixture.py} (the fixture), the C++ and Rust parity tests (\texttt{make test-code}),
\texttt{ml\_infer.accuracy()}, and, once, \texttt{bench\_infer.py} under \texttt{nice}.
\end{tutsteps}
\textbf{What to change next.} Store the forest's nodes in breadth-first order to improve locality, or evaluate several trees
at once with vector instructions (Book~13, chapter~14), and measure.
\end{tutorial}

\section{Build: the inference kernels}\label{bld:ml:low-latency-inference:mlinfer}

\begin{build}
\bfield{Purpose} Models in the fast path with the Python reference's exact results.
\bfield{Interface} Python: \texttt{export\_forest}, \texttt{forest\_predict}, \texttt{write\_forest}, \texttt{read\_forest},
\texttt{cpp\_branches}, \texttt{MLP}, \texttt{fit\_mlp(model, X, y, epochs, seed, qat, lr, batch, bits)}, \texttt{quantise(model,
X\_calib, bits)}, \texttt{int8\_forward}, \texttt{write\_mlp}; \texttt{make\_mlinfer\_fixture.py}. C++20: \texttt{mlinfer::Forest},
\texttt{read\_forest}, \texttt{Mlp}, \texttt{read\_mlp}, the generated \texttt{forest\_branches}. Rust: \texttt{Forest},
\texttt{Mlp}, \texttt{requant}.
\bfield{Rules} No allocation, exception or virtual call on the prediction path; parity on shared vectors is exact; timings are
measured, labelled with the machine, and never asserted.
\bfield{Acceptance tests} \texttt{code/firm/mlinfer/tests/} (Python: export equals LightGBM, fixture reproducible, integer
reference against the float model), \texttt{cpp/mlinfer\_test.cpp} and \texttt{rust} (\texttt{cargo test}): exact parity of the
flat forest, the generated branches and the int8 network on the 200 vectors.
\bfield{Stretch} SIMD evaluation of several trees; per-channel \omterm{def:ml:low-latency-inference:quant}{quantisation}; a hardware-description export of the int8 network.
\end{build}

\begin{omsources}
\item B.~Jacob and co-authors, ``Quantization and training of neural networks for efficient integer-arithmetic-only
inference'', arXiv:1712.05877, 2018.
\item M.~Nagel and co-authors, ``A white paper on neural network quantization'', arXiv:2106.08295, 2021.
\item N.~Asadi, J.~Lin and A.~P.~de Vries, ``Runtime optimizations for tree-based machine learning models'', \emph{IEEE
Transactions on Knowledge and Data Engineering}, 2014.
\item C.~Lucchese and co-authors, ``QuickScorer: a fast algorithm to rank documents with additive ensembles of regression
trees'', \emph{SIGIR}, 2015.
\item J.~Duarte and co-authors, ``Fast inference of deep neural networks in FPGAs for particle physics'', arXiv:1804.06913,
2018.
\end{omsources}

\section{Exercises}

\begin{exercise}[$\star$]\label{exo:ml:low-latency-inference:1}
A symmetric int8 tensor has scale $s = \max|w|/127$. What is the largest rounding error of one weight, as a share of $\max|w|$?
\end{exercise}

\begin{exercise}[$\star$]\label{exo:ml:low-latency-inference:2}
From \cref{fig:ml:inf:latency}, what is the ratio of the generated C++ forest's 99th-percentile latency to the library call's?
\end{exercise}

\begin{exercise}[$\star$]\label{exo:ml:low-latency-inference:3}
Why can the compiled forest reproduce LightGBM's prediction exactly while a compiled float network usually cannot reproduce
PyTorch's?
\end{exercise}

\begin{exercise}[$\star\star$]\label{exo:ml:low-latency-inference:4}
Write the fixed-point multiplier and shift for a real multiplier of 0.0123. Why use integers for it at all?
\end{exercise}

\begin{exercise}[$\star\star$]\label{exo:ml:low-latency-inference:5}
Why did \omterm{def:ml:low-latency-inference:quant}{quantisation-aware training} lose a little at eight bits and win at four?
\end{exercise}

\begin{exercise}[$\star\star$]\label{exo:ml:low-latency-inference:6}
\emph{Find the flaw.} ``Our compiled model's median latency is 800 nanoseconds, measured by timing 10\,000 predictions in a
batch and dividing by 10\,000.''
\end{exercise}

\begin{exercise}[$\star\star\star$]\label{exo:ml:low-latency-inference:7}
\emph{Coding.} Quantise the network to four bits after training with one weight scale per output unit instead of one per
tensor. How much of the loss against the float model does it recover?
\end{exercise}

\begin{exercise}[$\star\star\star$]\label{exo:ml:low-latency-inference:8}
With symmetric scales $s_x$, $s_w$ and $s_y$ for a layer's input, weights and output, write the layer's float output in terms of
the integers $X_i$, $W_i$ and an integer bias $B$, and derive the scale at which $B$ must be stored and the multiplier that maps
the accumulator to the output's integers.
\end{exercise}

\section{Problem: Two Microseconds for a Forest}

\begin{problem}[{Weekend problem --- the model is not the latency}]\label{pb:ml:low-latency-inference:1}
The chapter's forest, network, kernels and measurements.

\medskip\noindent\textbf{Part I --- The models.}
\begin{enumerate}
\item What are the two models, and how good are they?
\item How is the forest exported, and in how many forms?
\item Why are the kernels exact?
\item How large is each model in memory?
\end{enumerate}

\medskip\noindent\textbf{Part II --- Latency.}
\begin{enumerate}[resume]
\item Give the median and 99th-percentile latencies of each path.
\item Where does the library call's time go?
\item Why do the generated branches beat the loop?
\item What does the Rust version cost against C++, and why?
\end{enumerate}

\medskip\noindent\textbf{Part III --- \omterm{def:ml:low-latency-inference:quant}{Quantisation} and batching.}
\begin{enumerate}[resume]
\item What does int8 cost in information coefficient, after training and with \omterm{def:ml:low-latency-inference:quant}{quantisation-aware training}?
\item What changes at four bits?
\item How does LightGBM's latency scale with batch size?
\item When would you batch?
\end{enumerate}

\medskip\noindent\textbf{Part IV --- The verdict.}
\begin{enumerate}[resume]
\item State the \emph{named result}: the ratio of the compiled model's tail latency to the library call's, and the IC lost to
int8 \omterm{def:ml:low-latency-inference:quant}{quantisation} with and without \omterm{def:ml:low-latency-inference:quant}{quantisation-aware training}.
\item Which model would you deploy in a five-microsecond budget, and in what form?
\item What must the registry hold for a compiled model?
\item How would you make these measurements trustworthy on a production machine?
\item What would change with 5\,000 trees?
\item Where would programmable hardware be worth its cost?
\item How would a feature-store skew (chapter~24) show up at this layer?
\item In one sentence: where does inference time go?
\end{enumerate}
\end{problem}

\section{Interview questions}

\begin{interviewq}[$\star$ \iqroles{developer, mle}]\label{iq:ml:low-latency-inference:1}
How would you make a gradient-boosted model predict in a few microseconds?
\end{interviewq}

\begin{interviewq}[$\star\star$ \iqroles{mle}]\label{iq:ml:low-latency-inference:2}
Explain int8 \omterm{def:ml:low-latency-inference:quant}{quantisation}: scales, zero points, accumulators and requantisation.
\end{interviewq}

\begin{interviewq}[$\star\star$ \iqroles{developer}]\label{iq:ml:low-latency-inference:3}
What rules would you impose on code in the inference hot path?
\end{interviewq}

\begin{interviewq}[$\star\star$ \iqroles{mle, developer}]\label{iq:ml:low-latency-inference:4}
How do you test that a compiled model is the model you trained?
\end{interviewq}

\begin{interviewq}[$\star\star$ \iqroles{developer}]\label{iq:ml:low-latency-inference:5}
When does batching help \omterm{def:ml:low-latency-inference:latency}{inference latency}, and when does it hurt?
\end{interviewq}

\begin{interviewq}[$\star\star\star$ \iqroles{mle}]\label{iq:ml:low-latency-inference:6}
\omterm{def:ml:low-latency-inference:quant}{Post-training quantisation} or \omterm{def:ml:low-latency-inference:quant}{quantisation-aware training}? How would you decide?
\end{interviewq}
